\section{The problem and main results}

The kissing-number problem asks for the largest number $K(n)$ of unit
balls that can touch a central unit ball in Euclidean $n$-space without
overlapping interiors. Dividing the outer centres by two gives an equivalent
problem on the unit sphere: place as many points as possible at angular
distance at least $60^\circ$:
\[
K(n)=\max\{|X|:X\subset S^{n-1},\
\langle x,y\rangle\le\tfrac12\quad(x\ne y)\}.
\]
The equivalence follows from $|x-y|^2=2-2\langle x,y\rangle$. It connects
local sphere packing to spherical codes, and makes lattices, error-correcting
codes and finite geometry natural sources of constructions
\cite{conway1999sphere,dgs,leech-sloane1971}.

Higher dimension allows more directions, but the number of pair constraints
grows quadratically with the size of a configuration. An additional point
may obstruct many old points, so searching for an isolated gap often fails
to improve a large construction. Lattices and codes organize these
constraints: sign patterns on a support share one distance argument, and
vectors in a shell share norm and product restrictions. The problem becomes
how to exploit that organization without freezing every choice it contains.
We retain structures controlling most pairs while changing finite parts
that determine the cardinality.

\subsection{Results}

\begin{theorem}\label{thm:main}
For every dimension $n$ in Table~\ref{tab:results}, there is a kissing configuration of $N_n$ points, where $N_n$ is the lower bound in the second column. Thus $K(n)\ge N_n$.
\end{theorem}

The constructions in this report give the following lower bounds. The comparison column uses the fixed public sources recorded on 24–25 September 2026.

\begin{longtable}{rrrrl}
\caption{Explicit lower bounds and their public comparison values.}\label{tab:results}\\
\toprule
Dimension & Lower bound & Public comparison & Increase & Source\\\midrule
\endfirsthead
\toprule
Dimension & Lower bound & Public comparison & Increase & Source\\\midrule
\endhead
\bottomrule\endfoot
32 & 347584 & 346432 & 1152 & \cite{brouwer} \\
33 & 363968 & 362048 & 1920 & \cite{brouwer} \\
34 & 384196 & 381124 & 3072 & \cite{brouwer} \\
35 & 409676 & 409548 & 128 & \cite{brouwer} \\
36 & 484760 & 484568 & 192 & \cite{brouwer} \\
37 & 498024 & 496232 & 1792 & \cite{brouwer} \\
38 & 591900 & 591612 & 288 & \cite{kissingnumbers} \\
39 & 763668 & 756116 & 7552 & \cite{kissingnumbers} \\
43 & 2553792 & 2545056 & 8736 & \cite{sunwang-sections,crosssections} \\
45 & 7380090 & 7379838 & 252 & \cite{sunwang-sections,crosssections} \\
\end{longtable}


\subsection{Background on the constructions}

The first route starts with the three-layer construction of Edel, Rains
and Sloane~\cite{ers}. Its count combines a dense binary code, a family of
weight-eight supports and a root system. The Cheng--Sloane code supplies
the length-32 dense layer~\cite{chengsloane}. Constant-weight constructions
and comparisons are developed in the work of Brouwer, Shearer, Sloane and
Smith and in the subsequent tables~\cite{bsss,brouwer}. Echols recently
improved the bounds in dimensions 32, 33, 34 and 37 by support search
\cite{echols}. Two further choices remain: can candidates share the cost
of deleting old supports, and must a complete sign bundle remain
indivisible? The local replacement principle has a direct precursor in
the free zones of Takhanov--Yun~\cite[Section V]{signed-johnson}.
The constructions here combine that principle with paired Hadamard
transfer and explicit ERS embeddings. The seven code-based results
give the support families, transverse selections and boundary conditions
needed to turn these choices into larger configurations.

The second route uses the minimal shell of the Leech lattice. A lifting
combines a direction in the mother space with a low-dimensional tail;
the two products together determine compatibility. The lifting framework
of Cohn, Jiao, Kumar and Torquato~\cite[Theorem 7.5]{cjkt} and subsequent
constructions show how directions can be reused. Our dimension-38 result
starts from Kravatskiy's public 591612-point configuration
\cite{kissingnumbers}. All minimal Leech lines already occur in its
lifted layer, but its equator is empty. Instead of reusing those same
minimal directions again, the search moves to another shell and seeks
directions compatible with every lifted point. This adds 288 points.

The third route takes a small section of an extremal 48-dimensional
lattice and studies its perpendicular shell or the projections of selected
fibres. Conway--Sloane sections and the antipode construction establish
this connection~\cite{laminated,antipode}; the exact counts of Sun and Wang
supply the public comparisons in dimensions 43 and 45
\cite{sunwang-sections,crosssections}. The mother shell has 52416000
vectors. The later analysis of embedded $E_8$ sections by
Dorofeev, Sun and Wang~\cite{dorofeev} provides the relevant parent
geometry for the 43-dimensional calculation. The large shell makes
direct enumeration expensive. Moreover, a section's
abstract root type and determinant do not determine its embedding.
Spherical-design theory~\cite{dgs,nebe-designs} supplies exact low-degree
moment identities constraining projection multiplicities. We use them
to determine a 43-dimensional section count and to turn a small explicit
witness set into a lower bound of millions of points in dimension 45.

\subsection{Main contributions and construction principles}

The research gives explicit lower bounds in \ResultCount{} dimensions and
mathematical methods linking construction searches to exact proofs.
In layered codes, shared conflicts determine the net gain of a joint
support exchange. Paired Hadamard transfer turns weight-four codes into
compatible local weight-eight signed codes, allowing support selection
and sign assignments to change together. The 256-point replacement in
dimension 35 attains the capacity bound of its eleven-coordinate model;
the 576-point replacement in dimension 36 has a direct four-colour
construction. Dimension 37 combines support enlargement and signed
replacement, with gains of 1664 and 128 points. Each construction includes
an explicit finite object and a proof of its boundary compatibility.

In dimension 38, the key is to identify geometric space still available
outside the lifted layer. We select 144 compatible lines from another
Leech shell and prove both their mutual angular separation and their
compatibility with every lifted point. This preserves the base configuration
and adds 288 equatorial points. The construction resolves the choice of
directions, the inner-product constraints between shells and the common
coordinate realization of both layers.

Dimensions 43 and 45 turn geometric construction problems into exact
counting problems. In dimension 43, an embedded $E_8$ parent constrains
the fibre multiplicities of a $D_5$ child; rational moment identities
determine a count of 2,553,792 points. In dimension 45, a rational
inequality for the target fibres is combined with a common-neighbour
graph for witness search. The resulting 298 explicit vectors establish
a lower bound of 7,380,090. These constructions exhibit two distinct
uses of counting: determining a child's exact size from its parent
structure, and using a small constructible set to control the size
of a large projected configuration.

Qiushi Engine autonomously conducted the construction searches, mathematical
analysis and computational verification. Its autonomous experimental research
is described in \cite{qiushi-optics}. After the result
table, the report proceeds through code constructions, equatorial completion
and section counts. Each chapter introduces the object being changed,
derives its compatibility conditions, and explains how the finite data
specify the resulting construction.
